\begin{derivation}[从\eqnrefs{eq:ssb-general-potential-invariance-dp68}到\eqnrefs{eq:ssb-general-hessian-zero-mode-dp68}]
把一组实标量场坐标记为 $\phi_i$。连续整体对称性的无穷小作用写成

\begin{equation*}
 \delta\phi_i=\epsilon^a(T^a)_{ij}\phi_j,
\end{equation*}
其中 $\epsilon^a$ 是无穷小实参数，$T^a$ 是该实场表示中的生成元。势能在整个场空间上保持不变，因此对任意场构型都有
\begin{align*}
 0=\delta V
 &=\frac{\partial V}{\partial\phi_i}\,\delta\phi_i,\\
 &=\epsilon^a\frac{\partial V}{\partial\phi_i}(T^a)_{ij}\phi_j.
\end{align*}
由于各 $\epsilon^a$ 相互独立，每个生成元分别满足函数恒等式
\begin{equation*}
 \frac{\partial V}{\partial\phi_i}(T^a)_{ij}\phi_j=0.
\end{equation*}
对 $\phi_k$ 再求一次偏导，乘积法则给
\begin{equation*}
 0=
 \frac{\partial^2V}{\partial\phi_k\partial\phi_i}
 (T^a)_{ij}\phi_j
 +\frac{\partial V}{\partial\phi_i}(T^a)_{ik}.
\end{equation*}
令 $v_i$ 为平稳真空，即
\begin{equation*}
 \left.\frac{\partial V}{\partial\phi_i}\right|_{\phi=v}=0.
\end{equation*}
第二项在真空处消失，于是
\begin{align*}
 0
 &=\left.\frac{\partial^2V}{\partial\phi_k\partial\phi_i}
 \right|_{v}(T^av)_i,\\
 &=H_{ki}(v)(T^av)_i,
\end{align*}
其中
\begin{equation*}
 H_{ki}(v)\equiv
 \left.\frac{\partial^2V}{\partial\phi_k\partial\phi_i}\right|_v
\end{equation*}
是势能在真空处的 Hessian 矩阵。若生成元 $T^a$ 被真空自发破缺，则 $T^av\neq0$；上式说明非零向量 $T^av$ 属于 $H(v)$ 的零空间。对具有标准正定动能项的实标量场，二次拉氏量中的质量平方矩阵正比于该 Hessian 矩阵，所以这个真空轨道切向方向对应零质量模。Goldstone 零模因此来自“势能沿连续对称轨道保持常数”这一一般几何事实，而不依赖墨西哥帽势的具体形状。
\end{derivation}

\begin{derivation}[从\eqnrefs{eq:general-gauge-boson-mass-matrix-dp11}到\eqnrefs{eq:gauge-mass-positive-quadratic-form-dp68}]

正文\eqnrefs{eq:general-gauge-boson-mass-matrix-dp11} 给出生成元空间中的质量矩阵。要判断它是否可能产生负质量平方，不需要先对角化；只需把它与任意实向量 \(x_a\) 收缩，并把结果改写成真空向量的范数平方。


定义
\begin{equation*}
 (M_g^2)_{ab}
 =\frac{g_{\rm G}^2}{\hbar c}
 v_R^\dagger\{T^a,T^b\}v_R.
\end{equation*}
对任意实伴随空间向量 $x_a$，令
\begin{equation*}
 X\equiv x_aT^a.
\end{equation*}
若采用厄米生成元 $T^{a\dagger}=T^a$，则 $X^\dagger=X$。质量矩阵对应的二次型逐步化为
\begin{align*}
 x_a(M_g^2)_{ab}x_b
 &=\frac{g_{\rm G}^2}{\hbar c}
 v_R^\dagger x_ax_b\{T^a,T^b\}v_R,\\
 &=\frac{g_{\rm G}^2}{\hbar c}
 v_R^\dagger\{X,X\}v_R,\\
 &=\frac{2g_{\rm G}^2}{\hbar c}v_R^\dagger X^2v_R,\\
 &=\frac{2g_{\rm G}^2}{\hbar c}(Xv_R)^\dagger(Xv_R),\\
 &=\frac{2g_{\rm G}^2}{\hbar c}\|Xv_R\|^2\ge0.
\end{align*}
因此 $M_g^2$ 是实对称正半定矩阵，树级谱不会从这一机制产生负质量平方的矢量模。

等号成立当且仅当
\begin{equation*}
 Xv_R=0.
\end{equation*}
于是质量矩阵的零空间恰由所有湮灭真空的 Lie 代数方向组成，也就是未破缺子代数。反之，若 $Xv_R\neq0$，对应二次型严格为正，该生成元方向属于破缺子空间，并为相应规范玻色子提供非零质量。破缺方向上的 Goldstone 切向自由度与规范场纵向偏振重组为有质量矢量表示；这一线性代数结论不依赖某个具体规范群的显式矩阵对角化。
\end{derivation}
